\chapter[Lezione 19 --- 17 novembre 2026]{Lezione 19: parità, asintoti e limiti notevoli}
\label{chap:lezione-19}

\section{Parità e limiti unilateri}
\label{sec:lez19-parita-limiti}

La parità introdotta nella \cref{sec:lez18-parita} permette
di collegare i limiti unilateri nell'origine.

\begin{proposition}[Limiti unilateri di una funzione pari]
\label{prop:lez19-limiti-funzione-pari}
Sia \(a>0\) e si ponga
\[
  A:=(-a,a)\setminus\{0\},
\]
e sia \(f\colon A\longrightarrow\R\) pari.
Per ogni \(\ell\in\overline{\R}\), si ha
\[
  \lim_{x\to0^-}f(x)=\ell
  \quad\Longleftrightarrow\quad
  \lim_{x\to0^+}f(x)=\ell.
\]
\end{proposition}

\begin{proof}
Si supponga che
\[
  \lim_{x\to0^-}f(x)=\ell.
\]
Per \(x\to0^+\), posto \(u:=-x\), si ha \(u\to0^-\).
Poiché \(f\) è pari, per la
\cref{def:lez18-funzione-pari} risulta
\[
  f(x)=f(-u)=f(u).
\]
Applicando il \cref{thm:lez16-limite-sostituzione},
si ottiene
\[
  \lim_{x\to0^+}f(x)
  =\lim_{u\to0^-}f(u)=\ell.
\]
L'implicazione inversa segue con lo stesso argomento,
scambiando i limiti unilateri.
\end{proof}

\begin{proposition}[Limiti unilateri di una funzione dispari]
\label{prop:lez19-limiti-funzione-dispari}
Sia \(a>0\) e si ponga
\[
  A:=(-a,a)\setminus\{0\},
\]
e sia \(f\colon A\longrightarrow\R\) dispari.
Per ogni \(\ell\in\overline{\R}\), con la convenzione
\(-(+\infty)=-\infty\) e \(-(-\infty)=+\infty\), vale
\[
  \lim_{x\to0^-}f(x)=\ell
  \quad\Longleftrightarrow\quad
  \lim_{x\to0^+}f(x)=-\ell.
\]
\end{proposition}

\begin{proof}
Si supponga che
\[
  \lim_{x\to0^-}f(x)=\ell.
\]
Per \(x\to0^+\), posto \(u:=-x\), si ha \(u\to0^-\).
Poiché \(f\) è dispari, per la
\cref{def:lez18-funzione-dispari} risulta
\[
  f(x)=f(-u)=-f(u).
\]
Applicando il \cref{thm:lez16-limite-sostituzione}
e il \cref{thm:lez14-limite-prodotto},
si ottiene
\[
  \lim_{x\to0^+}f(x)
  =-\lim_{u\to0^-}f(u)=-\ell.
\]
L'implicazione inversa segue con lo stesso argomento,
scambiando i limiti unilateri.
\end{proof}

Per il \cref{thm:lez14-limite-unilateri},
se \(f\) è pari ed esiste uno dei due limiti unilateri,
esiste anche il limite bilatero con lo stesso valore.
Se \(f\) è dispari ed entrambi i limiti unilateri sono
finiti, il limite bilatero esiste se e solo se sono
entrambi nulli.

\section{Parità e coefficienti degli asintoti}
\label{sec:lez19-parita-asintoti}

La simmetria rispetto all'asse delle ordinate o
all'origine consente di confrontare le rette
asintotiche ai due estremi del dominio.

\begin{exercise}[Coefficienti asintotici di funzioni pari e dispari]
\label{ex:lez19-coefficienti-asintotici}
Sia \(A\subseteq\R\) non limitato superiormente e
inferiormente, simmetrico rispetto a zero, e sia
\(f\colon A\longrightarrow\R\) pari oppure dispari.
Si supponga che esistano finiti i limiti
\[
  m_+:=\lim_{x\to+\infty}\frac{f(x)}x,
  \qquad
  m_-:=\lim_{x\to-\infty}\frac{f(x)}x,
\]
\[
  q_+:=\lim_{x\to+\infty}\bigl(f(x)-m_+x\bigr),
  \qquad
  q_-:=\lim_{x\to-\infty}\bigl(f(x)-m_-x\bigr),
\]
con i limiti intesi per \(x\in A\).
Determinare le relazioni tra \(m_+,m_-\) e tra
\(q_+,q_-\), distinguendo i due casi.
\end{exercise}

\begin{proof}[Soluzione dell'esercizio]
\emph{Primo caso: \(f\) pari.}
Ponendo \(u:=-x\), per \(x\to-\infty\) si ha
\(u\to+\infty\). Dalla parità segue
\[
\begin{aligned}
  m_-
  &=\lim_{u\to+\infty}\frac{f(-u)}{-u}
   =-\lim_{u\to+\infty}\frac{f(u)}u
   =-m_+,\\
  q_-
  &=\lim_{u\to+\infty}\bigl(f(-u)+m_+(-u)\bigr)
   =\lim_{u\to+\infty}\bigl(f(u)-m_+u\bigr)
   =q_+.
\end{aligned}
\]

\emph{Secondo caso: \(f\) dispari.}
Ponendo \(u:=-x\), per \(x\to-\infty\) si ha
\(u\to+\infty\). Dalla disparità segue
\[
\begin{aligned}
  m_-
  &=\lim_{u\to+\infty}\frac{f(-u)}{-u}
   =-\lim_{u\to+\infty}\frac{-f(u)}u
   =m_+,\\
  q_-
  &=\lim_{u\to+\infty}\bigl(f(-u)+m_+u\bigr)
   =-\lim_{u\to+\infty}\bigl(f(u)-m_+u\bigr)
   =-q_+.
\end{aligned}
\]
Pertanto, per una funzione pari,
\[
  m_-=-m_+ ,\qquad q_-=q_+,
\]
mentre, per una funzione dispari,
\[
  m_-=m_+ ,\qquad q_-=-q_+.
\]
\end{proof}

Quando i limiti indicati esistono, le rette
\(y=m_+x+q_+\) e \(y=m_-x+q_-\) descrivono
il comportamento asintotico della funzione,
rispettivamente per \(x\to+\infty\) e \(x\to-\infty\).

\clearpage

\section{Limiti notevoli derivati dal numero di Nepero}
\label{sec:lez19-limiti-nepero}

Si dimostrano alcuni dei limiti enunciati nel
\cref{thm:lez17-limiti-notevoli}, a partire dal
\cref{thm:lez17-numero-nepero}.

Per le potenze a base variabile positiva ed esponente
reale si utilizza l'identità
\[
  f(x)^{g(x)}=e^{g(x)\log f(x)},
\]
valida nel dominio
\[
  D=(D_f\cap D_g)\cap\{x\in\R:f(x)>0\}.
\]

\begin{proposition}[Limite di Nepero per \(x\to\pm\infty\)]
\label{prop:lez19-nepero-infinito}
Si ha
\[
  \lim_{x\to+\infty}\left(1+\frac1x\right)^x=e,
  \qquad
  \lim_{x\to-\infty}\left(1+\frac1x\right)^x=e.
\]
Il dominio della funzione è
\[
  D=(-\infty,-1)\cup(0,+\infty).
\]
\end{proposition}

\begin{proof}
La base deve essere positiva e \(x\ne0\); pertanto
\[
  1+\frac1x>0
  \quad\Longleftrightarrow\quad
  x\in(-\infty,-1)\cup(0,+\infty).
\]

\emph{Primo caso: \(x\to+\infty\).}
Per \(x>1\), si ponga \(k:=\lfloor x\rfloor\).
Per la \cref{def:floor},
\[
  1\le k\le x<k+1,
\]
da cui
\[
  1+\frac1{k+1}<1+\frac1x\le1+\frac1k.
\]
Poiché le basi sono maggiori di \(1\), si ottiene
\[
  \left(1+\frac1{k+1}\right)^k
  \le\left(1+\frac1{k+1}\right)^x
  <\left(1+\frac1x\right)^x
  \le\left(1+\frac1k\right)^x
  <\left(1+\frac1k\right)^{k+1}.
\]
Per \(x\to+\infty\) si ha \(k\to+\infty\).
Osservando che
\[
\begin{aligned}
  \left(1+\frac1{k+1}\right)^k
    &=\frac{\left(1+\frac1{k+1}\right)^{k+1}}
            {1+\frac1{k+1}},\\[4pt]
  \left(1+\frac1k\right)^{k+1}
    &=\left(1+\frac1k\right)^k
      \left(1+\frac1k\right),
\end{aligned}
\]
per il \cref{thm:lez17-numero-nepero} e i
\cref{thm:lez14-limite-prodotto,thm:lez14-limite-rapporto}
i due estremi tendono a \(e\).
Dal \cref{thm:lez12-tre-funzioni} segue
\[
  \lim_{x\to+\infty}\left(1+\frac1x\right)^x=e.
\]

\emph{Secondo caso: \(x\to-\infty\).}
Posto \(t:=-x\), si ha \(t\to+\infty\) e, per \(t>1\),
\[
\begin{aligned}
  \left(1+\frac1x\right)^x
  &=\left(1-\frac1t\right)^{-t}
   =\left(\frac{t}{t-1}\right)^t\\
  &=\left(1+\frac1{t-1}\right)^{t-1}
    \left(1+\frac1{t-1}\right).
\end{aligned}
\]
Il primo fattore tende a \(e\) per il limite
appena dimostrato, mentre il secondo tende a \(1\).
Per il \cref{thm:lez14-limite-prodotto},
si conclude
\[
  \lim_{x\to-\infty}\left(1+\frac1x\right)^x=e.
\]
\end{proof}

\begin{proposition}[Limite di Nepero in zero]
\label{prop:lez19-nepero-zero}
Si ha
\[
  \lim_{x\to0}(1+x)^{1/x}=e.
\]
\end{proposition}

\begin{proof}
Si distinguono i limiti unilateri.
Per \(x\to0^+\), posto \(t:=1/x\), si ha
\(t\to+\infty\); quindi, per la
\cref{prop:lez19-nepero-infinito},
\[
  \lim_{x\to0^+}(1+x)^{1/x}
  =\lim_{t\to+\infty}\left(1+\frac1t\right)^t=e.
\]
Per \(x\to0^-\), si ha invece \(t\to-\infty\) e
\[
  \lim_{x\to0^-}(1+x)^{1/x}
  =\lim_{t\to-\infty}\left(1+\frac1t\right)^t=e.
\]
I due limiti coincidono; per il
\cref{thm:lez14-limite-unilateri}, segue la tesi.
\end{proof}

\begin{exercise}[Limiti esponenziali con parametro]
\label{ex:lez19-limiti-parametro}
Dimostrare mediante sostituzione che, per ogni
\(\alpha\in\R\),
\begin{enumerate}
  \item
  \[
    \lim_{n\to+\infty}
      \left(1+\frac\alpha n\right)^n=e^\alpha;
  \]
  \item
  \[
    \lim_{x\to+\infty}
      \left(1+\frac\alpha x\right)^x=e^\alpha;
  \]
  \item
  \[
    \lim_{x\to-\infty}
      \left(1+\frac\alpha x\right)^x=e^\alpha.
  \]
\end{enumerate}
Le espressioni sono considerate nei rispettivi domini,
nei quali la base è positiva.
\end{exercise}


\begin{proof}[Soluzione dell'esercizio]
Se \(\alpha=0\), tutte le espressioni sono
identicamente uguali a \(1=e^0\).

Sia dunque \(\alpha\ne0\). Per ogni \(u>0\), si ha
\[
  u^\alpha=e^{\alpha\log u}.
\]
Poiché \(\lim_{u\to e}\log u=1\) per
l'\cref{ex:lez18-limiti-logaritmo}, dal
\cref{thm:lez18-limite-esponenziale} e dai teoremi
sulle operazioni dei limiti segue
\[
  \lim_{u\to e}u^\alpha=e^\alpha.
\]

\emph{Primo e secondo punto.}
Il primo limite è la restrizione del secondo ai
valori naturali \(x=n\). È quindi sufficiente dimostrare
il secondo limite.

Ponendo \(t:=x/\alpha\), per \(x\to+\infty\)
si ha \(t\to+\infty\) se \(\alpha>0\),
mentre \(t\to-\infty\) se \(\alpha<0\).
Poiché \(x=\alpha t\), segue
\[
  \left(1+\frac\alpha x\right)^x
  =\left[\left(1+\frac1t\right)^t\right]^\alpha.
\]
Per la \cref{prop:lez19-nepero-infinito},
l'espressione tra parentesi quadre tende a \(e\)
in entrambi i casi. Pertanto
\[
  \lim_{x\to+\infty}
    \left(1+\frac\alpha x\right)^x=e^\alpha.
\]
Restringendo \(x\) ai valori \(n\in\N\), per il
\cref{thm:lez17-teorema-ponte} si ottiene anche
\[
  \lim_{n\to+\infty}
    \left(1+\frac\alpha n\right)^n=e^\alpha.
\]

\emph{Terzo punto.}
Ponendo nuovamente \(t:=x/\alpha\), per
\(x\to-\infty\) si ha \(t\to-\infty\) se
\(\alpha>0\), mentre \(t\to+\infty\) se
\(\alpha<0\). Poiché \(x=\alpha t\), segue
\[
  \left(1+\frac\alpha x\right)^x
  =\left[\left(1+\frac1t\right)^t\right]^\alpha.
\]
Per la \cref{prop:lez19-nepero-infinito},
l'espressione tra parentesi quadre tende a \(e\)
in entrambi i casi. Pertanto
\[
  \lim_{x\to-\infty}
    \left(1+\frac\alpha x\right)^x=e^\alpha.
\]
\end{proof}

\section{Limite notevole delle potenze in zero}
\label{sec:lez19-limite-potenze}

Si dimostra ora il sesto punto del
\cref{thm:lez17-limiti-notevoli}, utilizzando
il limite del logaritmo già stabilito nello stesso teorema.

\begin{proposition}[Limite notevole delle potenze]
\label{prop:lez19-limite-potenze}
Per ogni \(\alpha\in\R\), si ha
\[
  \lim_{x\to0}\frac{(1+x)^\alpha-1}{x}=\alpha.
\]
\end{proposition}

\begin{proof}
Si considera la potenza reale \((1+x)^\alpha\)
nel dominio \(x>-1\).

\emph{Primo caso: \(\alpha=0\).}
Il numeratore è nullo per ogni \(x>-1\),
perciò il limite vale \(0=\alpha\).

\emph{Secondo caso: \(\alpha\ne0\).}
Si ponga
\[
  t:=(1+x)^\alpha-1.
\]
Poiché
\[
  (1+x)^\alpha=e^{\alpha\log(1+x)},
\]
per i limiti del logaritmo
(\cref{ex:lez18-limiti-logaritmo})
e dell'esponenziale (\cref{thm:lez18-limite-esponenziale})
si ha \(t\to0\) per \(x\to0\).
Inoltre, per \(x\ne0\) sufficientemente vicino a \(0\),
si ha \(t\ne0\) e
\[
  \log(1+t)=\alpha\log(1+x).
\]
Pertanto
\[
\begin{aligned}
  \frac{(1+x)^\alpha-1}{x}
  &=\frac{t}{\log(1+t)}
    \cdot\frac{\alpha\log(1+x)}{x}.
\end{aligned}
\]
Dal quarto punto del \cref{thm:lez17-limiti-notevoli},
applicato alla base \(e\), segue
\[
  \lim_{z\to0}\frac{\log(1+z)}z=1,
  \qquad
  \lim_{z\to0}\frac{z}{\log(1+z)}=1.
\]
Applicando il \cref{thm:lez16-limite-sostituzione}
alla funzione \(t=(1+x)^\alpha-1\) e il
\cref{thm:lez14-limite-prodotto}, si conclude
\[
  \lim_{x\to0}\frac{(1+x)^\alpha-1}{x}
  =1\cdot\alpha=\alpha.
\]
\end{proof}
